% Part: normal-modal-logic
% Chapter: tableaux
% Section: introduction

\documentclass[../../../include/open-logic-section]{subfiles}

\begin{document}

\olfileid{nml}{tab}{int}

\olsection{પરિચય}

!!^{tableau}s એ !!{signed
  formula}sનાં ચોક્કસ પ્રકારનાં (નીચે તરફ શાખા પાડતાં) વૃક્ષો છે.
ચિહ્નિત સૂત્ર એટલે સત્યમૂલ્ય-ચિહ્ન ($\True$ અથવા
$\False$) અને !!a{sentence}ની જોડી:
\[
\sFmla{\True}{!A} \text{ અથવા } \sFmla{\False}{!A}.
\]
!!^a{tableau} કેટલીક \emph{ધારણાઓ}થી શરૂ થાય છે. પછીનું દરેક
!!{signed formula} અનુમાનનો કોઈ નિયમ લાગુ કરીને મળે છે.
કેટલાક નિયમો વૃક્ષની શાખાના છેડે એક અથવા વધુ
!!{signed formula}s ઉમેરે છે; બીજા નિયમો બે નવા છેડા ઉમેરી
બે શાખાઓ બનાવે છે. કોઈ નિયમથી મળતા !!{signed formula}sમાંનું
!!{formula} જે !!{signed formula} પર નિયમ લાગુ થયો તેના
સૂત્ર કરતાં ઓછું સંકુલ હોય છે. કોઈ શાખામાં
$\sFmla{\True}{!A}$ અને $\sFmla{\False}{!A}$ બંને હોય,
તો તે શાખાને \emph{બંધ} કહીએ છીએ. !!a{tableau}ની દરેક શાખા
બંધ હોય તો આખું !!{tableau} બંધ છે. બંધ !!{tableau} એક
!!a{derivation} છે, જે બતાવે છે કે !!{tableau}ની શરૂઆતમાં લીધેલા
!!{signed formula}sનો ગણ અસંતોષ્ય છે. એના આધારે
$\Proves$ સંબંધ વ્યાખ્યાયિત કરી શકાય:
$\Gamma \Proves !A$ જો અને માત્ર જો એવો કોઈ સાંત ગણ
~$\Gamma_0 = \{!B_1,
\dots, !B_n\} \subseteq \Gamma$ હોય કે નીચેની ધારણાઓ માટે
બંધ !!{tableau} મળે:
\[
\{\sFmla{\False}{!A}, \sFmla{\True}{!B_1}, \dots, \sFmla{\True}{!B_n}\}.
\]

નિદર્શ્ય તર્કશાસ્ત્રો માટે !!{signed
formula}ની સંકલ્પના વિસ્તૃત કરવી પડે છે અને
~\iftag{prvBox}{$\Box$\iftag{prvDiamond}{ અને
    $\Diamond$}{}}{$\Diamond$}ને આવરી લેતા નિયમો ઉમેરવા પડે છે.
\sourcecorrection{OLMD-051}{મૂળમાં અંદરના
\texttt{\string\iftag\{prvDiamond\}} માટે જરૂરી ત્રીજું, ખાલી
શાખા-આર્ગ્યુમેન્ટ ખૂટે છે; પસંદગી-મેક્રો યોગ્ય રીતે ચાલે તે માટે
તે ઉમેર્યું છે.}
ચિહ્ન ($\True$ અથવા $\False$) ઉપરાંત, મોડલ !!{tableau}sનાં
!!{formula}sને \emph{ઉપસર્ગો}~$\sigma$ પણ હોય છે. ઉપસર્ગો ધન
પૂર્ણાંકોની ખાલી ન હોય તેવી શ્રેણીઓ છે, એટલે
$\sigma \in (\PosInt)^* \setminus
\{\emptyseq\}$. આવા ઉપસર્ગો લખતી વખતે આપણે આસપાસના
$\tuple{\ }$ છોડીએ છીએ અને જુદા !!{element}s વચ્ચે
$,$ને બદલે~$.$ મૂકીએ છીએ. જો $\sigma$ ઉપસર્ગ હોય, તો
$\sigma.n$ એ $\sigma
\concat \tuple{n}$ છે; દાખલા તરીકે, $\sigma = 1.2.1$ હોય
તો $\sigma.3$ એ $1.2.1.3$ છે. આમ, દાખલા તરીકે,
\[
\sFmla{\True}{\Box !A \lif !A}[1.2]
\]
એ \emph{ઉપસર્ગવાળું !!{signed formula}} છે (અથવા ટૂંકમાં
\emph{ઉપસર્ગવાળું !!{formula}}).

સહજ રીતે, ઉપસર્ગ એ નિદર્શનમાં એવા જગતનું નામ આપે છે જે
!!a{tableau}ની શાખા પરનાં !!{formula}sને સંતોષી શકે; અને જો
$\sigma$ કોઈ જગતનું નામ હોય, તો $\sigma.n$ એ~$\sigma$ નામના
જગતમાંથી પ્રાપ્ય જગતનું નામ આપે છે.

\end{document}
